\documentclass[../Main.tex]{subfiles}
\begin{document}

\section{Installation and Setup Guide}

This appendix provides detailed instructions for installing, configuring, and operating the ShieldX system in a production environment.

\subsection{Prerequisites}

Before deploying ShieldX, ensure the following prerequisites are met:

\begin{itemize}
    \item \textbf{Agent (Linux)}: Ubuntu 20.04+ or any Linux distribution with Python 3.10+, root/sudo access for packet capture and nftables management, and network connectivity to the ShieldX-web backend.
    \item \textbf{Backend (Linux/Windows)}: Python 3.10+, MySQL 8.0 database (can be Dockerized), and network access for agents to connect.
    \item \textbf{Frontend}: Node.js 18+ and npm for building the React dashboard. The built static files can be served by any web server (Nginx, Apache, etc.).
\end{itemize}

\subsection{Agent Installation}

\subsubsection{Step 1: Clone the Repository}

\begin{verbatim}
git clone <repository-url> ShieldX
cd ShieldX
\end{verbatim}

\subsubsection{Step 2: Set Up Python Environment}

\begin{verbatim}
python3 -m venv .venv
source .venv/bin/activate
pip install -r requirements.txt
\end{verbatim}

\subsubsection{Step 3: Configure the Agent}

Copy the example configuration file and edit it with your backend details:

\begin{verbatim}
cp agent_config.yaml.example agent_config.yaml
\end{verbatim}

The configuration file supports the following parameters:

\begin{itemize}
    \item \texttt{backend\_url}: The URL of the ShieldX-web backend API (e.g., http://192.168.1.100:8000).
    \item \texttt{agent\_id}: A unique identifier for this agent (auto-generated if not specified).
    \item \texttt{interface}: The network interface to capture packets on (e.g., eth0).
    \item \texttt{sniff\_duration}: Duration of each capture round in seconds (default: 120).
    \item \texttt{log\_level}: Logging verbosity (DEBUG, INFO, WARNING, ERROR).
\end{itemize}

\subsubsection{Step 4: Place the ML Model}

Ensure the XGBoost model file is in the project root directory:

\begin{verbatim}
cp xgboost_model.json ShieldX/
\end{verbatim}

The agent expects the model file at the path specified in \texttt{src/domains/pipeline.py} (default: \texttt{xgboost\_model.json} in the current working directory).

\subsubsection{Step 5: Run the Agent}

For full orchestrator mode (recommended):

\begin{verbatim}
sudo python -m src.main
\end{verbatim}

For lightweight scheduling mode (detection only, no firewall management):

\begin{verbatim}
sudo python -m src.agent_scheduler
\end{verbatim}

The agent requires root privileges for two operations: (i) raw socket access for packet capture via Scapy, and (ii) nftables rule management for domain whitelist enforcement.

\subsection{Backend Installation}

\subsubsection{Step 1: Set Up MySQL Database}

Using Docker (recommended):

\begin{verbatim}
docker run --name shieldx-mysql -e MYSQL_ROOT_PASSWORD=root \
  -e MYSQL_DATABASE=shieldx -p 3306:3306 -d mysql:8.0
\end{verbatim}

\subsubsection{Step 2: Configure and Start the Backend}

\begin{verbatim}
cd services
cp .env.example .env
# Edit .env with your MySQL connection details
pip install -r requirements.txt
uvicorn main:app --host 0.0.0.0 --port 8000
\end{verbatim}

The backend will automatically create the required database tables on first startup using SQLModel's ORM.

\subsection{Frontend Deployment}

\subsubsection{Step 1: Install Dependencies and Build}

\begin{verbatim}
cd shieldx-web-frontend
npm install
echo "VITE_API_URL=http://localhost:8000" > .env
npm run build
\end{verbatim}

\subsubsection{Step 2: Serve the Built Files}

The built files in the \texttt{dist/} directory can be served by any static web server. Example Nginx configuration:

\begin{verbatim}
server {
    listen 80;
    server_name shieldx.example.com;
    root /path/to/shieldx-web-frontend/dist;
    index index.html;
    location /api/ {
        proxy_pass http://localhost:8000;
    }
}
\end{verbatim}

\subsection{Operation and Maintenance}

\subsubsection{Agent Logging}

The agent logs to both stdout and a rotating file (\texttt{agent.log}). The log format includes timestamps, log levels, and module names:

\begin{verbatim}
2026-06-22 10:30:00,123 - pipeline - INFO - Starting capture round...
2026-06-22 10:32:00,456 - pipeline - INFO - Processed 15 flows
2026-06-22 10:32:00,789 - pipeline - WARNING - Malicious flow detected: 192.168.1.50:443
\end{verbatim}

\subsubsection{Monitoring Agent Health}

The agent sends heartbeat signals to the backend every 60 seconds. The heartbeat payload includes system information (hostname, IP address, OS version, CPU and memory usage) and the current agent status. Administrators can monitor agent health through the web dashboard's Agents page.

\subsubsection{Updating the ML Model}

To update the ML model without restarting the agent:

\begin{enumerate}
    \item Replace the \texttt{xgboost\_model.json} file with the new model.
    \item The agent automatically detects the file modification on the next capture round.
    \item If hot-reloading fails, restart the agent with \texttt{sudo python -m src.main}.
\end{enumerate}

\subsubsection{Troubleshooting Common Issues}

\begin{itemize}
    \item \textbf{"Permission denied" for packet capture}: Ensure the agent is run with root privileges or that the Python binary has CAP\_NET\_RAW capability: \texttt{sudo setcap cap\_net\_raw+ep \$(readlink -f \$(which python))}.
    \item \textbf{Backend connection refused}: Verify that the backend server is running and that the \texttt{backend\_url} in \texttt{agent\_config.yaml} is correct. Check for firewall rules blocking port 8000.
    \item \textbf{No flows detected}: Verify that the agent is monitoring the correct network interface and that there is actual TLSv1.3 traffic on port 443. Use \texttt{tcpdump -i eth0 port 443} to verify packet capture independently.
    \item \textbf{Memory usage growing}: The flow garbage collection parameters in \texttt{flow\_session.py} may need tuning. Reduce \texttt{EXPIRED\_UPDATE} interval or increase the flow timeout threshold.
    \item \textbf{Model loading fails}: Ensure the \texttt{xgboost\_model.json} file exists and is compatible with the installed XGBoost version. Re-download the model if it appears corrupted.
\end{itemize}

\subsection{Security Considerations for Production}

\begin{itemize}
    \item \textbf{HTTPS}: Configure TLS termination on the reverse proxy (Nginx) for the backend API to encrypt agent-backend communication.
    \item \textbf{Authentication}: Add JWT-based authentication to the backend API to prevent unauthorized access. The current implementation does not include authentication.
    \item \textbf{Database Credentials}: Use environment variables or a secrets manager for MySQL credentials rather than hardcoding them in configuration files.
    \item \textbf{Agent Isolation}: Run each agent in a dedicated system user account with minimal privileges. Use Linux namespaces or containers for additional isolation.
    \item \textbf{Log Retention}: Configure log rotation to prevent disk exhaustion. The default configuration keeps 30 days of logs before rotation.
\end{itemize}

\end{document}
